% TOP_PROBLEM: {"id": 30007280, "title": "Rationality of the canonical unramified Arthur filtration", "category_id": 1, "category": {"id": 1, "name": "number_theory", "display_name": "Number Theory", "description": "Properties of integers, prime numbers, Diophantine equations.", "slug": "number-theory", "order_index": 1, "created_at": "2026-07-31T15:26:25.670Z"}, "status": "open"}
% FIRST_POSTED: 2026-10-06
% ENTRY_KIND: statement_only
% !TeX program = lualatex
\documentclass[11pt]{article}
\usepackage[margin=25mm]{geometry}
\usepackage{fontspec}
\setmainfont{Latin Modern Roman}
\usepackage{amsmath,amssymb,mathtools}
\usepackage{unicode-math}
\setmathfont{Latin Modern Math}
\usepackage{xurl,hyperref}
\hypersetup{colorlinks=true,urlcolor=blue}
\setcounter{secnumdepth}{0}
\setlength{\parindent}{0pt}
\setlength{\parskip}{5pt}
\setlength{\emergencystretch}{3em}
\newcommand{\R}{\mathbb R}
\newcommand{\N}{\mathbb N}
\newcommand{\Z}{\mathbb Z}
\newcommand{\Q}{\mathbb Q}
\newcommand{\C}{\mathbb C}
\newcommand{\D}{\mathbb D}
\newcommand{\F}{\mathbb F}
\newcommand{\sref}[2]{\hyperlink{source:#1:#2}{[S#2]}}
\newcommand{\cataloguescope}[1]{#1}
\begin{document}
\section*{Rationality of the canonical unramified Arthur filtration}\label{30007280}
\subsection{Definitions and mathematical statement}
\paragraph{Statement recorded in the supplied source.}
Let X/Fq be a smooth projective geometrically connected curve and G/Fq split connected semisimple. Assume that the Lie algebra admits an invariant nondegenerate symmetric form nondegenerate on every Levi center, Chevalley restriction is an isomorphism for every Levi, semisimple Lie-element scheme centralizers are Levi subgroups, and every nilpotent Lie element over every extension belongs to the Lie algebra of the unipotent radical of an extension-defined parabolic. Let $A_E$ be finitely supported E-valued functions on the finite-field G-bundle isomorphism classes, $E=Qbar_\ell $. For any closed invariant subset Y of the dual Lie-algebra nilpotent cone, let $F_Y$ $A_E$ be the image of categorical Frobenius trace from the corresponding derived-Satake nilpotent-support subcategory. Must E $tensor_Q$ $(A_Q$ intersect $F_Y$ $A_E)$ map isomorphically to $F_Y$ $A_E$ for every ell different from p, yielding a common rational filtration independent of ell?
\subsection{Short English statement}
Let X/Fq be a smooth projective geometrically connected curve and G/Fq split connected semisimple. Assume that the Lie algebra admits an invariant nondegenerate symmetric form nondegenerate on every Levi center, Chevalley restriction is an isomorphism for every Levi, semisimple Lie-element scheme centralizers are Levi subgroups, and every nilpotent Lie element over every extension belongs to the Lie algebra of the unipotent radical of an extension-defined parabolic. Let $A_E$ be finitely supported E-valued functions on the finite-field G-bundle isomorphism classes, $E=Qbar_\ell $. For any closed invariant subset Y of the dual Lie-algebra nilpotent cone, let $F_Y$ $A_E$ be the image of categorical Frobenius trace from the corresponding derived-Satake nilpotent-support subcategory. Must E $tensor_Q$ $(A_Q$ intersect $F_Y$ $A_E)$ map isomorphically to $F_Y$ $A_E$ for every ell different from p, yielding a common rational filtration independent of ell?
\subsection{Sources}
\begin{itemize}
\item[S1] Original question source or mathematical context.
\url{https://www.claymath.org/library/monographs/cmim04c.pdf}.
\item[S2] Original question source or mathematical context.
\url{https://doi.org/10.4171/owr/2021/39}.
\item[S3] Original question source or mathematical context.
\url{https://doi.org/10.4171/owr/2024/54}.
\item[S4] OpenAI, manuscript citation page. Source of the formulation motivating this local problem.
\url{https://github.com/openai/math/blob/adc7f1241b42e322a6451854ab7e4b4c146bf78a/preprints/Rationality-of-the-Canonical-Unramified-Arthur-Filtration-September-24-2026/README.md}.
\end{itemize}
\end{document}
